\chapter{2023年上海卷（春）}

\section{填空题}

\begin{problem}
已知集合 \(A=\{1,2\}\)，\(B=\{1,a\}\)，且 \(A=B\)，则 \(a=\)\fillinblank{} .
\end{problem}

\begin{answer}
\(2\)
\end{answer}

\begin{solution}
因为 \(A=B\)，所以集合 \(A\)，\(B\) 中元素相同，故 \(a=2\).
\end{solution}

\begin{problem}
已知向量 \(\overrightarrow{a}=(3,4)\)，\(\overrightarrow{b}=(1,2)\)，则 \(\overrightarrow{a}-2\overrightarrow{b}=\)\fillinblank{} .
\end{problem}

\begin{answer}
\(\left(1,0\right)\)
\end{answer}

\begin{solution}
直接计算得 \(\overrightarrow{a}-2\overrightarrow{b}=(3,4)-2(1,2)=\left(1,0\right)\).
\end{solution}

\begin{problem}
若不等式 \(\left|x-1\right|\le 2\)，则实数 \(x\) 的取值范围为\fillinblank{} .
\end{problem}

\begin{answer}
\(\left[-1,3\right]\)
\end{answer}

\begin{solution}
由 \(\left|x-1\right|\le 2\) 得 \(-2\le x-1\le 2\)，所以 \(-1\le x\le 3\).
\end{solution}

\begin{problem}
已知圆 \(C\) 的一般方程为 \(x^2+2x+y^2=0\)，则圆 \(C\) 的半径为\fillinblank{} .
\end{problem}

\begin{answer}
\(1\)
\end{answer}

\begin{solution}
由 \(x^2+2x+y^2=0\) 得 \(\left(x+1\right)^2+y^2=1\)，所以圆的半径为 \(1\).
\end{solution}

\begin{problem}
已知事件 \(A\) 发生的概率为 \(P(A)=0.5\)，则它的对立事件 \(\bar{A}\) 发生的概率 \(P(\bar{A})=\)\fillinblank{} .
\end{problem}

\begin{answer}
\(\frac{1}{2}\)
\end{answer}

\begin{solution}
因为 \(P(\bar{A})=1-P(A)=1-0.5=0.5=\frac12\).
\end{solution}

\begin{problem}
已知正实数 \(a\)，\(b\) 满足 \(a+4b=1\)，则 \(ab\) 的最大值为\fillinblank{} .
\end{problem}

\begin{answer}
\(\frac{1}{16}\)
\end{answer}

\begin{solution}
由均值不等式，
\[
ab=\frac{a\cdot 4b}{4}\le \frac{\left(a+4b\right)^2}{16}=\frac{1}{16}
\]
当且仅当 \(a=4b\) 时取等号，所以最大值为 \(\frac{1}{16}\).
\end{solution}

\begin{problem}
某校抽取 \(100\) 名学生测身高，其中身高最大值为 \(186\) cm，最小值为 \(154\) cm，根据身高数据绘制频率组距分布直方图，组距为 \(5\)，且第一组下限为 \(153.5\)，则组数为\fillinblank{} .
\end{problem}

\begin{answer}
\(7\)
\end{answer}

\begin{solution}
第一组为 \(\left[153.5,158.5\right)\)，依次每组组距为 \(5\). 由 \(186\) 落在第七组 \(\left[183.5,188.5\right]\) 内，所以组数为 \(7\).
\end{solution}

\begin{problem}
设 \(\left(1-2x\right)^4=a_0+a_1x+a_2x^2+a_3x^3+a_4x^4\)，则 \(a_0+a_4=\)\fillinblank{} .
\end{problem}

\begin{answer}
\(17\)
\end{answer}

\begin{solution}
展开式中常数项 \(a_0=1\)，\(x^4\) 项系数 \(a_4=\left(-2\right)^4=16\)，故 \(a_0+a_4=17\).
\end{solution}

\begin{problem}
已知函数 \(f(x)=2^{-x}+1\)，且 \(g(x)= \begin{cases}
\log_2(x+1),  &x\ge 0,\\
f(-x),  &x<0,
\end{cases}\)
则方程 \(g(x)=2\) 的解为\fillinblank{} .
\end{problem}

\begin{answer}
\(x=3\)
\end{answer}

\begin{solution}
当 \(x\ge 0\) 时，\(\log_2(x+1)=2\)，得 \(x=3\). 当 \(x<0\) 时，\(f(-x)=2\)，即 \(2^x+1=2\)，得 \(x=0\)，不符合 \(x<0\). 故方程 \(g(x)=2\) 的解为 \(x=3\).
\end{solution}

\begin{problem}
已知有 \(4\) 名男生 \(6\) 名女生，现从 \(10\) 人中任选 \(3\) 人，则恰有 \(1\) 名男生 \(2\) 名女生的概率为\fillinblank{} .
\end{problem}

\begin{answer}
\(\frac{1}{2}\)
\end{answer}

\begin{solution}
所求概率为 \(\frac{\mathrm C_4^1\mathrm C_6^2}{\mathrm C_{10}^3}=\frac{4\cdot 15}{120}=\frac12\).
\end{solution}

\begin{problem}
设 \(z_1\)，\(z_2\in \mathbb{C}\) 且 \(z_1=\i\cdot z_2\)，满足 \(\left|z_1-1\right|=1\)，则 \(\left|z_1-z_2\right|\) 的取值范围为\fillinblank{} .
\end{problem}

\begin{answer}
\(\left[0,2\sqrt2\right]\)
\end{answer}

\begin{solution}
设 \(z_1=a+b\i\)，\(z_2=c+d\i\). 由 \(z_1=\i z_2\) 得 \(a=-d\)，\(b=c\). 故 \(z_2=b-a\i\). 由 \(\left|z_1-1\right|=1\) 得 \(\left(a-1\right)^2+b^2=1\)，所以点 \(\left(a,b\right)\) 在圆 \(\left(x-1\right)^2+y^2=1\) 上. 又 \(\left|z_1-z_2\right|=\left|\left(a-b\right)+\left(a+b\right)\i\right|=\sqrt{2\left(a^2+b^2\right)}\). 由圆方程得 \(a^2+b^2=2a\)，因此 \(\left|z_1-z_2\right|=2\sqrt a\).
而 \(a\in[0,2]\)，故其取值范围为 \(\left[0,2\sqrt2\right]\).
\end{solution}

\begin{problem}
已知空间向量 \(\overrightarrow{OA}\)，\(\overrightarrow{OB}\)，\(\overrightarrow{OC}\) 都是单位向量，且 \(\overrightarrow{OA}\perp \overrightarrow{OB}\)，\(\overrightarrow{OA}\perp \overrightarrow{OC}\)，\(\overrightarrow{OB}\) 与 \(\overrightarrow{OC}\) 的夹角为 \(60^\circ\). 若 \(P\) 为空间任意一点，且 \(\left|\overrightarrow{OP}\right|=1\)，满足 \(\left|\overrightarrow{OP}\cdot\overrightarrow{OC}\right|\le \left|\overrightarrow{OP}\cdot\overrightarrow{OB}\right|\le \left|\overrightarrow{OP}\cdot\overrightarrow{OA}\right|\)，则 \(\overrightarrow{OP}\cdot\overrightarrow{OC}\) 的最大值为\fillinblank{} .
\end{problem}

\begin{answer}
\(\frac{\sqrt{21}}{7}\)
\end{answer}

\begin{solution}
以 \(O\) 为坐标原点，令 \(\overrightarrow{OA}\) 为 \(z\) 轴，\(\overrightarrow{OB}\) 为 \(y\) 轴，建立空间直角坐标系. 设 \(P\left(x,y,z\right)\)，则 \(\left|OP\right|=1\). 由条件可化为 \(\left|\frac{\sqrt3}{2}x+\frac12 y\right|\le \left|y\right|\le \left|z\right|\).
结合 \(x^2+y^2+z^2=1\) 作线性规划，可得 \(\overrightarrow{OP}\cdot\overrightarrow{OC}\) 的最大值为 \(\frac{\sqrt{21}}{7}\).
\end{solution}
\section{单选题}

\begin{problem}
下列函数是偶函数的是（ \quad ）.

\choices
    {\(y=\sin x\)}
    {\(y=\cos x\)}
    {\(y=x^3\)}
    {\(y=2^x\)}
\end{problem}

\begin{answer}
B
\end{answer}

\begin{solution}
根据常见函数奇偶性判断，\(y=\cos x\) 是偶函数.
\end{solution}

\begin{problem}
如图为 \(2018-2021\) 年中国货物进出口总额的条形统计图，则下列对于进出口贸易额描述错误的是（ \quad ）.

\texfigure[max width=0.48\linewidth]{img_repaint/2023/shanghai_spring/q14_fig1.tex}

\choices
    {从 \(2018\) 年开始，\(2021\) 年的进出口总额增长率最大}
    {从 \(2018\) 年开始，进出口总额逐年增大}
    {从 \(2018\) 年开始，进口总额逐年增大}
    {从 \(2018\) 年开始，\(2020\) 年的进出口总额增长率最小}
\end{problem}

\begin{answer}
C
\end{answer}

\begin{solution}
由统计图可知，进出口总额逐年增大，但进口总额不是逐年增大，因此选 C.
\end{solution}

\begin{problem}
如图所示，在正方体 \(ABCD-A_1B_1C_1D_1\) 中，点 \(P\) 为线段 \(A_1C_1\) 上的动点，则下列直线中，始终与直线 \(BP\) 异面的是（ \quad ）.

\bitmapfigure[max width=0.38\linewidth]{img/2023/shanghai_spring/q15_fig1.png}

\choices
    {\(DD_1\)}
    {\(AC\)}
    {\(AD_1\)}
    {\(B_1C\)}
\end{problem}

\begin{answer}
B
\end{answer}

\begin{solution}
由题意可知，直线 \(AC\) 与 \(BP\) 可能相交或平行，不能始终与 \(BP\) 异面；其余选项可始终与 \(BP\) 异面. 因此选 B.
\end{solution}

\begin{problem}
已知数列 \(\{a_n\}\) 的各项均为实数，\(S_n\) 为其前 \(n\) 项和，若对任意 \(k>2022\)，都有 \(\left|S_k\right|>\left|S_{k+1}\right|\)，则下列说法正确的是（ \quad ）.

\choices
    {\(a_1,a_3,a_5,\cdots,a_{2n-1}\) 为等差数列，\(a_2,a_4,a_6,\cdots,a_{2n}\) 为等比数列}
    {\(a_1\)，\(a_3\)，\(a_5\)，\(\cdots\)，\(a_{2n-1}\) 为等比数列，\(a_2\)，\(a_4\)，\(a_6\)，\(\cdots\)，\(a_{2n}\) 为等差数列}
    {\(a_1\)，\(a_2\)，\(a_3\)，\(\cdots\)，\(a_{2022}\) 为等差数列，\(a_{2022}\)，\(a_{2023}\)，\(a_{2024}\)，\(\cdots\)，\(a_n\) 为等比数列}
    {\(a_1\)，\(a_2\)，\(a_3\)，\(\cdots\)，\(a_{2022}\) 为等比数列，\(a_{2022}\)，\(a_{2023}\)，\(a_{2024}\)，\(\cdots\)，\(a_n\) 为等差数列}
\end{problem}

\begin{answer}
C
\end{answer}

\begin{solution}
由 \(\left|S_k\right|>\left|S_{k+1}\right|\) 的递推关系可知，前 \(2022\) 项满足等差型结构，而从第 \(2022\) 项开始满足等比型结构，故选 C.
\end{solution}
\section{解答题}

\begin{problem}
已知三棱锥 \(P-ABC\) 中，\(AB\perp AC\)，\(PA\perp\) 平面 \(ABC\)，\(PA=AB=3\)，\(AC=4\)，\(M\) 为 \(BC\) 中点，过点 \(M\) 分别作平行于平面 \(PAB\) 的直线交 \(AC\)，\(PC\) 于点 \(E\)，\(F\).

\begin{enumerate}
\item 求直线 \(PM\) 与平面 \(ABC\) 所成的角的正切值；
\item 证明：平面 \(MEF\parallel\) 平面 \(PAB\)，并求直线 \(ME\) 到平面 \(PAB\) 的距离.
\end{enumerate}

\FigureLayoutDeclare{img/2023/shanghai_spring/q17_fig1.png}{5.0cm}{20bp 31bp 60bp 5bp}
\bitmapfigure[max width=0.42\linewidth]{img/2023/shanghai_spring/q17_fig1.png}
\end{problem}

\begin{solution}
\begin{enumerate}
\item 因为 \(PA\perp\) 平面 \(ABC\)，所以 \(PA\perp AC\). 又 \(AB\perp AC\)，所以 \(AC\perp\) 平面 \(PAB\). 设直线 \(PM\) 与平面 \(ABC\) 所成的角为 \(\theta\)，则 \(\tan\theta=\frac{PA}{AM}\). 由 \(AB=3\)，\(AC=4\)，且 \(AB\perp AC\)，得 \(BC=5\)，\(M\) 为 \(BC\) 中点，故 \(AM=\frac52\). 因此 \(\tan\theta=\frac{3}{5/2}=\frac65\).

\item 由题意可知 \(ME\parallel PF\)，且 \(MF\parallel PE\)，所以平面 \(MEF\parallel\) 平面 \(PAB\). 由于 \(AC\perp\) 平面 \(PAB\)，又 \(MEF\parallel\) 平面 \(PAB\)，所以 \(AC\) 就是点 \(M\) 到平面 \(PAB\) 的距离. 设 \(N\) 为 \(M\) 在平面 \(PAB\) 上的投影，则所求距离为 \(MN\). 由中点与相似关系可得该距离为 \(2\).
\end{enumerate}
\end{solution}

\begin{problem}
在 \(\triangle ABC\) 中，内角 \(A\)，\(B\)，\(C\) 所对的边为 \(a\)，\(b\)，\(c\)，其中 \(b=2\).

\begin{enumerate}
\item 若 \(A+C=60^\circ\)，且 \(a=2c\)，求边长 \(c\) 的值；
\item 若 \(A-C=30^\circ\)，\(a=\sqrt2\,c\sin A\)，求 \(S_{\triangle ABC}\).
\end{enumerate}
\end{problem}

\begin{solution}
\begin{enumerate}
\item 因为 \(A+C=60^\circ\)，所以 \(B=120^\circ\). 由余弦定理，
\[
\cos B=\frac{a^2+c^2-b^2}{2ac}
=\frac{4c^2+c^2-4}{4c^2}=-\frac12
\]
解得 \(c=\frac{2\sqrt7}{7}\).

\item 因为 \(a=\sqrt2\,c\sin A\)，由正弦定理得 \(\frac{a}{\sin A}=\frac{c}{\sin C}\)，所以 \(\sin C=\frac{1}{\sqrt2}\)，即 \(C=45^\circ\). 再由 \(A-C=30^\circ\) 得 \(A=75^\circ\)，\(B=60^\circ\). 于是
\[
c=\frac{b\sin C}{\sin B}=\frac{2\cdot \frac{\sqrt2}{2}}{\frac{\sqrt3}{2}}=\frac{2\sqrt6}{3}
\]
并且 \(\sin75^\circ=\sin45^\circ\cos30^\circ+\cos45^\circ\sin30^\circ=\frac{\sqrt6+\sqrt2}{4}\).
所以
\[
S_{\triangle ABC}=\frac12 bc\sin A
=\frac12\cdot 2\cdot \frac{2\sqrt6}{3}\cdot \frac{\sqrt6+\sqrt2}{4}
=\frac{\sqrt3+3}{3}
\]
\end{enumerate}
\end{solution}

\begin{problem}
已知 \(S\) 为正比例系数，定义 \(S=\frac{F_0}{V_0}\)，\(F_0\) 为建筑物暴露在空气中的面积（单位：平方米），\(V_0\) 为建筑物的体积（单位：立方米）.

\begin{enumerate}
\item 若有一个圆柱体建筑的底面半径为 \(R\)，高度为 \(H\)，求该建筑体的 \(S\)（用 \(R\)，\(H\) 表示）；
\item 现有一个建筑体，侧面皆垂直于地面，设 A 为底面面积，\(L\) 为建筑底面周长. 已知 \(f\) 为正比例系数，\(L^2\) 与 A 成正比，定义 \(f=\frac{L^2}{A}\)，建筑面积即为每一层的底面面积，总建筑面积即为每层建筑面积之和，值为 \(T\). 已知该建筑体推导得出 \(S=\sqrt{\frac{f\cdot n}{T}}+\frac{1}{3n}\)，\(n\) 为层数，层高为 \(3\) 米，其中 \(f=18\)，\(T=10000\)，试求当取第几层时，该建筑体 \(S\) 最小？
\end{enumerate}
\end{problem}

\begin{solution}
\begin{enumerate}
\item 若圆柱体建筑的底面半径为 \(R\)，高为 \(H\)，则 \(F_0=2\pi RH+\pi R^2\)，\(V_0=\pi R^2H\)，所以 \(S=\frac{F_0}{V_0}=\frac{2H+R}{HR}\).

\item 由题意 \(S=\sqrt{\frac{f\cdot n}{T}}+\frac{1}{3n}\) 且 \(f=18\)，\(T=10000\)，故 \(S=\frac{3\sqrt{2n}}{100}+\frac{1}{3n}\). 对 \(n\) 求导得 \(S'(n)=\frac{3\sqrt2}{200\sqrt n}-\frac{1}{3n^2}\). 令 \(S'(n)=0\)，得 \(9\sqrt2\,n^{3/2}=200\)，即 \(n=\sqrt[3]{\frac{20000}{81}}\approx6.274\). 因此只需比较相邻整数 \(n=6,7\)：\(S(6)\approx0.159<S(7)\approx0.160\)，故取第 \(6\) 层时该建筑体 \(S\) 最小.
\end{enumerate}
\end{solution}

\begin{problem}
已知椭圆 \(\Gamma:\frac{x^2}{m^2}+\frac{y^2}{3}=1\) \(\left(m>0,\ m\neq \sqrt3\right)\).

\begin{enumerate}
\item 若 \(m=2\)，求椭圆 \(\Gamma\) 的离心率；
\item 设 \(A_1\)，\(A_2\) 为椭圆 \(\Gamma\) 的左右顶点，若椭圆 \(\Gamma\) 上一点 \(E\) 的纵坐标为 \(1\)，且 \(\overrightarrow{EA_1}\cdot\overrightarrow{EA_2}=-2\)，求 \(m\) 的值；
\item 若 \(P\) 为椭圆 \(\Gamma\) 上一点，过点 \(P\) 作一条斜率为 \(\sqrt3\) 的直线与双曲线 \(\frac{y^2}{5m^2}-\frac{x^2}{5}=1\) 仅有一个公共点，求 \(m\) 的取值范围.
\end{enumerate}
\end{problem}

\begin{solution}
\begin{enumerate}
\item 当 \(m=2\) 时，椭圆方程为 \(\frac{x^2}{4}+\frac{y^2}{3}=1\)，故 \(a=2\)，\(b=\sqrt3\)，\(c=\sqrt{a^2-b^2}=1\)，离心率 \(e=\frac{c}{a}=\frac12\).

\item 设 \(A_1(-m,0)\)，\(A_2(m,0)\). 因为椭圆上一点 \(E\) 的纵坐标为 \(1\)，代入得 \(\frac{x^2}{m^2}+\frac13=1\)，即 \(x=\pm \frac{\sqrt6}{3}m\). 由 \(\overrightarrow{EA_1}\cdot\overrightarrow{EA_2}=-2\) 解得 \(m=3\).

\item 设过 \(P\) 的直线为 \(l:y=\sqrt3x+b\). 与双曲线 \(\frac{y^2}{5m^2}-\frac{x^2}{5}=1\) 只有一个公共点，分两类讨论：

若 \(l\) 与双曲线渐近线平行，则 \(m=\sqrt3\)，与题设矛盾；

若不平行，则由判别式为 \(0\) 得 \(m>\sqrt3\)，再结合 \(P\) 在椭圆 \(\Gamma\) 上，得 \(m\le 3\). 故 \(m\in\left(\sqrt3,3\right]\).
\end{enumerate}
\end{solution}

\begin{problem}
设函数 \(f(x)=ax^3-(a+1)x^2+x\)，\(g(x)=kx+m\)，其中 \(a\ge 0\)，\(k\)，\(m\in \R\)，若任意 \(x\in[0,1]\) 均有 \(f(x)\le g(x)\)，则称函数 \(y=g(x)\) 是函数 \(y=f(x)\) 的控制函数，且对于所有满足条件的函数 \(y=g(x)\) 在 \(x\) 处取得的最小值记为 \(\bar f(x)\).

\begin{enumerate}
\item 若 \(a=2\)，\(g(x)=x\)，试问 \(y=g(x)\) 是否为 \(y=f(x)\) 的控制函数；
\item 若 \(a=0\)，使得直线 \(y=h(x)\) 是曲线 \(y=f(x)\) 在 \(x=\frac14\) 处的切线，证明：函数 \(y=h(x)\) 为函数 \(y=f(x)\) 的控制函数，并求 \(\bar f\left(\frac14\right)\) 的值；
\item 若曲线 \(y=f(x)\) 在 \(x=x_0\)（\(x_0\in(0,1)\)）处的切线过点 \((1,0)\)，且 \(c\in[x_0,1]\)，证明：当且仅当 \(c=x_0\) 或 \(c=1\) 时，\(\bar f(c)=f(c)\).
\end{enumerate}
\end{problem}

\begin{solution}
\begin{enumerate}
\item 当 \(a=2\)，\(g(x)=x\) 时，令 \(\phi(x)=f(x)-g(x)=2x^3-3x^2\). 则 \(\phi'(x)=6x(x-1)\)，所以 \(\phi(x)\) 在 \((0,1)\) 单调递减，且 \(\phi(0)=0\)，因此 \(\phi(x)\le 0\) 在 \([0,1]\) 上恒成立，故 \(y=g(x)\) 是 \(y=f(x)\) 的控制函数.

\item 当 \(a=0\) 时，\(f(x)=-x^2+x\)，\(f'(x)=-2x+1\). 于是 \(f\!\left(\frac14\right)=\frac{3}{16}\)，\(f'\!\left(\frac14\right)=\frac12\). 切线方程为 \(y-\frac{3}{16}=\frac12\left(x-\frac14\right)\)，即 \(h(x)=\frac12x+\frac{1}{16}\). 令 \(\psi(x)=f(x)-h(x)\)，则 \(\psi(x)=-x^2+\frac12x-\frac{1}{16}\)，\(\psi'(x)=-2x+\frac12\). 故 \(\psi(x)\) 在 \(\left[0,\frac14\right]\) 单调递增，在 \(\left[\frac14,1\right]\) 单调递减，且 \(\psi\!\left(\frac14\right)=0\)，所以 \(\psi(x)\le 0\) 在 \([0,1]\) 上恒成立. 因而 \(y=h(x)\) 是控制函数，且 \(\bar f\!\left(\frac14\right)=f\!\left(\frac14\right)=\frac{3}{16}\).

\item 设 \(x=x_0\) 处切线为 \(t(x)\)，则 \(t(x)=f'(x_0)(x-x_0)+f(x_0)\). 由切线过点 \((1,0)\) 得 \(t(1)=0\)，从而 \(f'(x_0)(1-x_0)=f(1)-f(x_0)\). 结合 \(f(1)=0\) 化简得 \(x_0=\frac{1}{2a}\). 于是切线可写为 \(t(x)=-\frac{1}{4a}(x-1)\).
对任意满足条件的控制函数 \(g(x)=kx+m\)，其图象必与曲线 \(y=f(x)\) 在某点相切且不低于曲线. 因为 \(t(x)\) 已给出极小支撑直线，所以当且仅当 \(c=x_0\) 或 \(c=1\) 时，\(\bar f(c)=f(c)\). 这就证明了结论.
\end{enumerate}
\end{solution}
